\documentclass[11pt,a4paper]{article}
\usepackage[T1]{fontenc}
\usepackage[utf8]{inputenc}
\usepackage{lmodern}
\usepackage{amsmath,amssymb,amsthm,mathtools,mathrsfs}
\usepackage{graphicx,float,xcolor,multicol}
\usepackage[shortlabels]{enumitem}
\usepackage[margin=27mm]{geometry}
\usepackage{microtype,needspace,placeins}
\usepackage{tikz,tikz-cd}
\usetikzlibrary{arrows.meta,calc,positioning,patterns,decorations.pathreplacing}
\usepackage[colorlinks=true,linkcolor=teal,urlcolor=teal,citecolor=teal]{hyperref}
\setlength{\parindent}{0pt}
\setlength{\parskip}{.6em}
\setcounter{secnumdepth}{0}
\allowdisplaybreaks
\raggedbottom
\providecommand{\cross}{\times}
\DeclareUnicodeCharacter{2020}{\ensuremath{\dagger}}
\DeclareMathOperator{\supp}{supp}
\DeclareMathOperator*{\esssup}{ess\,sup}
\providecommand{\chapter}{\section}
\newenvironment{tool}[2]{\subsection*{Tool #1 --- #2}}{}
\newcommand{\ToolRef}[1]{Tool #1}
\newcommand{\FigureTag}[2]{\textit{Figure #1}}
\newcommand{\Result}[1]{\subsection*{#1}}
\newcommand{\Application}[1]{\subsubsection*{Application\if\relax\detokenize{#1}\relax\else\space--- #1\fi}}
\newcommand{\ProofHeading}{\par\textit{Proof.}\space}
\newcommand{\ProofOf}[1]{\par\textit{Proof of #1.}\space}
\newcommand{\RemarkHeading}{\par\textbf{Remark.}\space}
\newcommand{\NoteHeading}{\par\textbf{Note.}\space}
\newcommand{\ExerciseHeading}{\par\textbf{Exercise.}\space}

% Rendering support only; mathematical text is in each independent post.
\usepackage{booktabs,array,longtable}
\definecolor{HandoutBlue}{HTML}{2F6690}
\definecolor{HandoutNavy}{HTML}{17365D}
\newtheorem{theorem}{Theorem}
\newtheorem{lemma}[theorem]{Lemma}
\newtheorem{proposition}[theorem]{Proposition}
\newtheorem{corollary}[theorem]{Corollary}
\newtheorem{conjecture}[theorem]{Conjecture}
\newtheorem{definition}{Definition}
\newtheorem{example}{Example}
\newtheorem{namedtoolcounter}{Tool}
\newenvironment{namedtool}[1]{\begin{namedtoolcounter}[#1]}{\end{namedtoolcounter}}
\newtheorem*{claim}{Claim}
\newtheorem*{remark}{Remark}
\newtheorem*{exercise}{Exercise}
\newtheorem*{question}{Question}
\newtheorem*{observation}{Observation}
\newcommand{\SourceChapter}[2]{\section*{#2}}
\newcommand{\Lecture}[2]{\section*{Lecture #1: #2}}
\newcommand{\unclear}[1]{\textit{[unclear: #1]}}
\newcommand{\sourcegap}{\textit{[unfinished in the source]}}
\DeclareMathOperator{\dist}{dist}
\DeclareMathOperator{\id}{id}
\DeclareMathOperator{\Hol}{Hol}
\DeclareMathOperator{\Per}{Per}
\DeclareMathOperator{\Fix}{Fix}
\DeclareMathOperator{\diam}{diam}
\DeclareMathOperator{\Var}{Var}
\DeclareMathOperator{\Leb}{Leb}
\DeclareMathOperator{\Hom}{Hom}
\DeclareMathOperator{\End}{End}
\DeclareMathOperator{\Ann}{Ann}
\DeclareMathOperator{\Spec}{Spec}
\DeclareMathOperator{\Max}{Max}
\DeclareMathOperator{\Ob}{Ob}
\DeclareMathOperator{\Tor}{Tor}
\DeclareMathOperator{\Ass}{Ass}
\DeclareMathOperator{\Mor}{Mor}
\DeclareMathOperator{\Fun}{Fun}
\DeclareMathOperator{\Nat}{Nat}
\DeclareMathOperator{\Cone}{Cone}
\DeclareMathOperator{\MCG}{MCG}
\DeclareMathOperator{\Aut}{Aut}
\DeclareMathOperator{\Deck}{Deck}
\DeclareMathOperator{\SL}{SL}
\DeclareMathOperator{\PSL}{PSL}
\DeclareMathOperator{\Area}{Area}
\DeclareMathOperator{\im}{im}
\DeclareMathOperator{\PSh}{PSh}
\newcommand{\CRing}{\mathsf{CRings}}
\newcommand{\AMod}{A\text{-}\mathsf{Mod}}
\newcommand{\Ab}{\mathsf{Ab}}
\newcommand{\Z}{\mathbb Z}
\newcommand{\Q}{\mathbb Q}
\newcommand{\R}{\mathbb R}
\newcommand{\C}{\mathbb C}
\newcommand{\N}{\mathbb N}
\newcommand{\kk}{\mathbb k}
\renewcommand{\aa}{\mathfrak a}
\newcommand{\bb}{\mathfrak b}
\newcommand{\pp}{\mathfrak p}
\newcommand{\qq}{\mathfrak q}
\newcommand{\mm}{\mathfrak m}
\newcommand{\nn}{\mathfrak n}
\newcommand{\rad}{\mathrm r}
\newcommand{\cat}[1]{\mathsf{#1}}
\setcounter{secnumdepth}{0}
\allowdisplaybreaks

\graphicspath{{figures/ergodic-theory/}}
\begin{document}
\section*{Recurrence and Three Views of Ergodicity}
% BLOG-CONTENT-BEGIN
\subsection{Recurrence Theorems}
Of all the impositions on a measure-preserving system, the one that is perhaps
ill-motivated is the demand that $\mu(X)<\infty$. Let's look at why it matters.

\setcounter{theorem}{3}\begin{theorem}[Poincar\'e recurrence]
Let $(X,\mathcal B,\mu,T)$ be a measure-preserving system and let
$E\in\mathcal B$. Then there exists $F\subseteq E$, $F\in\mathcal B$, such that
$\mu(F)=\mu(E)$ and, for every $x\in F$, $T^n(x)\in E$ for infinitely many
$n\geq1$.
\end{theorem}
\begin{proof}
Let us first look at sets that escape $E$ in one go:
\[
  B:=\{x\in E:T^nx\notin E\text{ for every }n\geq1\}.
\]
Note that
\begin{align*}
  B&=E\cap T^{-1}(X\setminus E)\cap T^{-2}(X\setminus E)\cap\cdots,\\
  T^{-n}(B)&=T^{-n}(E)\cap T^{-n-1}(X\setminus E)
  \cap T^{-n-2}(X\setminus E)\cap\cdots.
\end{align*}
If $n>m$, then
\[
  T^{-m}(B)=\left(\bigcap_{k\geq1}T^{-m-k}(X\setminus E)\right)\cap T^{-m}(E),
\]
and there exists $k_0$
% Source page 2
such that $n=m+k_0$. Thus
$T^{-m}(B)\subseteq T^{-m-k_0}(X\setminus E)=T^{-n}(X\setminus E)$, and we
have $T^{-n}(B)\subseteq T^{-n}(E)$. Hence
\[
  T^{-n}(B)\cap T^{-m}(B)=\varnothing\qquad(n\neq m).
\]
As $T$ is $\mu$-invariant and $\mu(X)=1$, it follows that $\mu(B)=0$.
Consequently, for almost every $x\in E$, $T^{\circ n}(x)\in E$ for some $n$.
Thus there exists $F_1\subseteq E$ with $\mu(F_1)=\mu(E)$, where
\[
  F_1=\{x\in E:T^{\circ n}(x)\in E\text{ for some }n\}.
\]
One can similarly get $F_2,F_3,\ldots$, working with $T^2,T^3,\ldots$.
Since $\mu(E\setminus F_i)=0$ for every $i$, the set $F=\bigcap_i F_i\subseteq E$
satisfies $\mu(F)=\mu(E)$.

In $F$, for every $x$ and every $n$, there exists $k$ such that $T^{nk}(x)\in E$.
To prove that, for every $x\in F$, $T^nx\in E$ for infinitely many $n$, note
that $x\in F$ implies $x\in F_1$, so there is an $n_1$ such that
$T^{n_1}(x)\in E$. Also, $x\in F_{n_1+1}$ implies that there is an
$n_2>n_1$ such that $T^{n_2}(x)\in E$, and so on. Hence $T^nx\in E$ for
infinitely many $n$, for almost every $x\in E$.
\end{proof}

\begin{remark}
Notice that this implies the pigeonhole principle if $X=\{1,2,3,\ldots,n\}$,
$\mu$ is the uniform probability measure on $2^X$, and $T$ is any bijective
map. What we essentially did is lift the pigeonhole principle from the
category of sets to that of measure-preserving systems. Here, we do not call
such results PHP (pigeonhole principle); we call them recurrence theorems.
\end{remark}

% Source page 3
\begin{remark}
We crucially used the fact that $\mu$ is a bounded measure. This turns out to
be the right analogue of a compact metric-space structure on $X$ in a
topological dynamical system. Recurrence is an important feature, because
recurrence is the first emergence of a pattern in a measure-preserving
system. Thus $\mu$ being a probability measure ensures that no information is
lost under the transformation $T$ of $X$.
\end{remark}

\begin{remark}
Define $T\colon\mathbb R^2\to\mathbb R^2$ by $T(x):=x+1$. Then $T$ is
$m$-invariant, where $m$ is Lebesgue measure on $\mathbb R^2$. But notice that
the set $T(\mathbb D)$ escapes to infinity. In such cases we have nothing in
$\mathbb R$ to compare $\mathbb D$ with: $\mathbb D$ runs off to infinity.
But if you had a measure-preserving system, $\mathbb D$ would return to itself,
almost, and hence studying the orbit amounts to studying a finite set of sets
that have almost-everywhere overlap, providing us with a means to compare!
\end{remark}
% Author check: R^2, x+1, R, and the assertions about a finite collection and
% overlap are the manuscript's notation and claims, retained without repair.

\begin{remark}
Bounded measure need not mean that the set $X$ is bounded, as shall be
illustrated now.
\end{remark}

\setcounter{theorem}{4}\begin{lemma}[Another litmus test for measure-preserving systems]
$(X,T,\mu)$ is a measure-preserving system if and only if
\[
  \int f\,d\mu=\int f\circ T\,d\mu
  \qquad\text{for every }f\in L^\infty(\mu).
\]
\end{lemma}

% Source page 4
\begin{proof}
Let $f=\chi_B$. Then
\[
  \mu(B)=\int\chi_B\,d\mu=\int\chi_B(T)\,d\mu
  =\int\chi_{T^{-1}(B)}\,d\mu=\mu(T^{-1}(B)),
\]
so $T$ is $\mu$-invariant. Conversely, as $\mu(X)<\infty$,
$L^1(\mu)\supset L^2(\mu)\supset\cdots\supset L^\infty(\mu)$, and we have
$\int f\,d\mu=\int f\circ T\,d\mu$ for all simple functions $f$.
Now let $f\in L^1(\mu)$. There exist $f_n\nearrow f$, so the dominated
convergence theorem gives
\[
  \int f\,d\mu=\int f\circ T\,d\mu
  \qquad\text{for every }f\in L^\infty(\mu).
\]
\end{proof}
% Author check: approximation by f_n increasing to f, and invocation of
% dominated convergence for a general L^1 function, are retained as written.

Now define the following measure on $\mathbb R$:
\[
  \mu([a,b]):=\int_a^b\frac{dx}{\pi(1+x^2)}.
\]
Let $(\mathbb R/\mathbb Z,m_{S^1})$ be the usual $S^1$ with Lebesgue measure.
The map $T\colon\mathbb R\to\mathbb R$ given by
\[
  T(x)=\frac12\left(x-\frac1x\right),\qquad T(0)=0,
\]
is $\mu$-invariant, because, for $f\in L^1(\mu)$, we have
\[
  \int_{-\infty}^{\infty}f(T(x))\frac{dx}{\pi(1+x^2)}
  =\int_{-\infty}^{\infty}f(T(x))\frac{dT(x)}{\pi(1+T(x)^2)}.
\]
Set $y=T(x)$ to obtain $\int f\circ T\,d\mu=\int f\,d\mu$. It is now easy to
check that
\[
  \phi\colon(\mathbb R,\mathcal B_{\mathbb R},\mu,T)
  \longrightarrow(\mathbb R/\mathbb Z,\mathcal B_{S^1},m_{S^1},T_2),
  \qquad \phi(x)=\frac1\pi\arctan(x)+\frac12,
\]
is an isomorphism.
% Author check: the displayed change-of-variables equality and the claimed
% arctangent conjugacy are retained exactly in mathematical substance.
But $\mathbb R$ is by no means bounded. Thus a measure-preserving system
ensures that transformations do not blow points off to infinity, and that is
why it can accommodate unbounded sets in its realm as well.

% Source page 5
\subsection{Ergodicity}
As far as the behaviour of orbits goes, we know that if $E$ is measurable in
$X$, then, for almost every $x\in E$, $\{T^{\circ n}(x)\}_n$ returns to $E$
infinitely often. Let us extend this madness further. Say $x\in X$ and
$E\subset X$ is measurable: can we have some sort of recurrence for the orbit
of $x$, $\{T^{\circ n}(x)\}_n$, if $x\notin E$?

There are some set-theoretic obstructions to such a recurrence, because if
$T^{-1}(E)=E$, the orbit of $x$ shall not intersect $E$. So let us replace
$X$ by $E$ and $E$ by some measurable $F\subseteq E$. Notice immediately that
we are looking for some sort of irreducibility criterion here.

\setcounter{definition}{4}\begin{definition}[Ergodicity]
Let $(X,\mathcal B,\mu,T)$ be a measure-preserving system. We say $T$ is
\emph{ergodic} if, for every $B\in\mathcal B$ such that $T^{-1}(B)=B$, we have
$\mu(B)=0$ or $\mu(B)=1$.
\end{definition}

We shall now try to understand ergodicity from three essentially distinct
perspectives, each revealing an interconnection.

\Needspace{14\baselineskip}
\subsubsection{Measure-Theoretic Ergodicity}

% Source page 6
\setcounter{theorem}{5}\begin{proposition}\label{prop:ergodicity-measure}
Let $(X,\mathcal B,\mu,T)$ be a measure-preserving system. The following are
equivalent:
\begin{enumerate}[label=(\arabic*)]
  \setcounter{enumi}{0}\item $T$ is ergodic.
  \setcounter{enumi}{1}\item For $B\in\mathcal B$, $\mu(T^{-1}(B)\mathbin\triangle B)=0$ implies
  $\mu(B)=0$ or $\mu(B)=1$ (almost ergodic).
  \setcounter{enumi}{2}\item $\mu(A)>0$ implies
  $\mu\bigl(\bigcup_{n=1}^{\infty}T^{-n}(A)\bigr)=1$
  (recurrence theorem).
  \setcounter{enumi}{3}\item $\mu(A)\mu(B)>0$ implies that there exists $n\geq1$ such that
  $\mu(T^{-n}A\cap B)>0$ (mixing criterion).
\end{enumerate}
\end{proposition}
\begin{proof}
\emph{$(1)\Rightarrow(2)$.} Set
\[
  C=\bigcap_{N=0}^{\infty}\bigcup_{n=N}^{\infty}T^{-n}(B).
\]
The idea is to pick $B$ such that $T^{-1}(B)=B$ almost everywhere, i.e.
$\mu(T^{-1}(B)\mathbin\triangle B)=0$, and construct a $C$ out of $B$ such that
$T^{-1}C=C$ and $\mu(C)=\mu(B)$. Note that
\[
  B\mathbin\triangle\left(\bigcup_{n=N}^{\infty}T^{-n}(B)\right)
  \subseteq\bigcup_{n=N}^{\infty}\bigl(B\mathbin\triangle T^{-n}(B)\bigr).
\]
Now
\begin{align*}
  \mu(T^{-1}(B)\mathbin\triangle B)
  &=\mu\bigl(T^{-1}(T^{-1}(B)\mathbin\triangle B)\bigr)\\
  &=\mu(T^{-2}(B)\mathbin\triangle T^{-1}(B))\\
  &=\cdots=\mu(T^{-i}B\mathbin\triangle T^{-i-1}(B))=0.
\end{align*}
Let $x\in B\mathbin\triangle T^{-n}(B)$. Then either
$x\in B\setminus T^{-n}(B)$ or $x\in T^{-n}(B)\setminus B$.
If $x\in B\setminus T^{-n}(B)$ and $x\in T^{-i}(B)$ for every
$0\leq i\leq k$, then $k\leq n$ and
% Source page 7
$x\in T^{-k}(B)\setminus T^{-k-1}(B)$. One can argue similarly for
$x\in T^{-n}(B)\setminus B$ to get
\[
  B\mathbin\triangle T^{-n}(B)
  \subseteq\bigcup_{i=0}^{n-1}\bigl(T^{-i}(B)\mathbin\triangle T^{-i-1}(B)\bigr),
\]
so $\mu(B\mathbin\triangle T^{-n}(B))=0$.
Set $C_N=\bigcup_{n=N}^{\infty}T^{-n}(B)$. Then, of course,
\[
  C_0\supseteq C_1\supseteq C_2\supseteq C_3\supseteq C_4\supseteq\cdots,
  \qquad \mu(C_N\mathbin\triangle B)=0.
\]
Since $C_N\searrow C$, we have $\mu(C\mathbin\triangle B)=0$, and hence
$\mu(C)=\mu(B)$. Also,
\[
  T^{-1}(C)=\bigcap_{N=0}^{\infty}\bigcup_{n=N}^{\infty}T^{-n-1}(B)
  =\bigcap_{N\geq1}C_N=\bigcap_{N\geq0}C_N=C.
\]
Therefore $\mu(C)=\mu(B)\in\{0,1\}$, by the ergodicity of $T$.

\emph{$(2)\Rightarrow(3)$.} Let $\mu(A)>0$ and set
$B=\bigcup_{n=1}^{\infty}T^{-n}(A)$. Then $T^{-1}(B)\subseteq B$.
Since $\mu$ is $T$-invariant, $\mu(T^{-1}(B)\mathbin\triangle B)=0$.
As $T^{-1}(A)\subseteq B$, we have $\mu(B)\neq0$, because
$\mu(T^{-1}(A))=\mu(A)>0$. Thus
\[
  \mu\left(\bigcup_{n=1}^{\infty}T^{-n}(A)\right)=1.
\]
Therefore, any point $x\in X$ recurs on a given set of positive measure via
an ergodic map.
% Author check: the source says "any point", rather than "almost every point".

\emph{$(3)\Rightarrow(4)$.} If $\mu(A)\mu(B)>0$, then $\mu(A)>0$ and
$\mu(B)>0$, so $\mu\bigl(\bigcup_{n=1}^{\infty}T^{-n}(A)\bigr)=1$.

% Source page 8
Since the $T^{-n}(A)$ cover $X$ almost everywhere, we have
\[
  0<\mu(B)<\mu\left(B\cap\bigcup T^{-n}(A)\right)
  \leq\sum_{n=1}^{\infty}\mu(B\cap T^{-n}(A)).
\]
Hence there exists $n\geq1$ such that $\mu(B\cap T^{-n}(A))>0$.
We shall see the relevance of such statements when we introduce the notion
of mixing systems.
% Author check: the strict inequality after mu(B) is present in the source.

\emph{$(4)\Rightarrow(1)$.} If $T^{-1}(A)=A$, then
\[
  0=\mu(A\cap(X\setminus A))
  =\mu(T^{-n}(A)\cap(X\setminus A))\qquad(n\geq1).
\]
Hence $\mu(A)=0$ or $\mu(X\setminus A)=0$, so $\mu(A)\in\{0,1\}$ and $T$ is
ergodic.
\end{proof}

\subsubsection{Operator-Theoretic Ergodicity}
Let us begin by associating an operator $U_T$ to $T$ as follows:
\[
  U_T\colon L^2(\mu)\longrightarrow L^2(\mu),\qquad U_T(f):=f\circ T.
\]
Now $L^2(\mu)$ is a Hilbert space with inner product $\langle\cdot,\cdot\rangle$.
Let $f_1,f_2\in L^2(\mu)$. Then
\[
  \langle U_Tf_1,U_Tf_2\rangle
  =\int(f_1\circ T)\overline{(f_2\circ T)}\,d\mu
  =\int f_1\overline{f_2}\,d\mu=\langle f_1,f_2\rangle,
\]
by the second litmus test. From this we see
$\langle f_1,U_T^*U_Tf_2\rangle=\langle f_1,f_1\rangle$, so
\[
  \langle f_1,(U_T^*U_T-\id)f_2\rangle=0
  \qquad\text{for all }f_1,f_2\in L^2(\mu).
\]
% Author check: the preceding right-hand inner product reads <f_1,f_1> in
% the manuscript, although the next deduction uses <f_1,f_2>.
A standard argument
% Source page 9
would tell you that
\[
  \sup_{\substack{\|f_1\|\leq1\\\|f_2\|\leq1}}
  \left|\langle f_1,(U_T^*U_T-\id)f_2\rangle\right|
  =\|U_T^*U_T-\id\|=0.
\]
Thus $U_T$ is a unitary operator, which goes by the name \emph{Koopman operator}.
% Author check: the source concludes "unitary" from U_T^* U_T = id, without
% an invertibility or surjectivity assumption on T.

\setcounter{theorem}{6}\begin{proposition}\label{prop:ergodicity-operator}
$T$ is ergodic if and only if the eigenvectors of $U_T$ associated to $1$ are
constants in $L^2_\mu(X)$.
\end{proposition}
\begin{proof}
If $B$ is such that $\mu(T^{-1}(B)\mathbin\triangle B)=0$, then
$\chi_B=\chi_B\circ T$ almost everywhere. Thus $U_T\chi_B=\chi_B$, so $\chi_B$
is a constant. Hence $\mu(B)=0$ or $\mu(B)=1$, and $T$ is ergodic.
If $U_Tf=f$ implies $f\equiv0$ almost everywhere, then $T$ is ergodic.
% Author check: the manuscript says f identically 0 here, not "f constant".

Conversely, let $f\colon X\to\mathbb R$ be measurable. Let $k\in\mathbb Z$
and $n\geq1$, and define
\[
  A_n^k:=\left\{x:f(x)\in\left[\frac{k}{n},\frac{k+1}{n}\right)\right\}.
\]
Now we know that
\[
  T^{-1}(A_n^k)\mathbin\triangle A_n^k
  \subseteq\{x:f\circ T(x)\neq f(x)\},
\]
which has measure zero if $U_Tf=f$. The ergodicity of $T$ implies
$\mu(A_n^k)\in\{0,1\}$, and we have $X=\bigsqcup_{k\in\mathbb Z}A_n^k$.
Thus there exists $k(n)\in\mathbb Z$ such that $\mu(A_n^{k(n)})=1$.
For every $n$,
\[
  f(x)\in\left[\frac{k(n)}{n},\frac{k(n)+1}{n}\right)
  \quad\text{almost everywhere}.
\]
Letting $n\to\infty$ shows that $f$ is a constant almost everywhere.
\end{proof}

\begin{remark}
Thus ergodicity can be rephrased in operator theory too. Properties possessing
such translations are called \emph{unitary properties}.
\end{remark}

% Source page 10
\subsubsection{Fourier-Analytic Ergodicity}
Here we impose a compact abelian group structure on $X$. Then it has a unique
translation-invariant measure $m_X$, called the \emph{Haar measure} of $X$.
We will not prove the existence of $m_X$ here. We shall further assume that,
given $f\in L^2_{m_X}(X)$, we have the following character-theoretic expansion:
\[
  f(x)=\sum_{\chi\in\widehat X}c_\chi\chi(x),
  \qquad \sum_{\chi\in\widehat X}|c_\chi|^2=\|f\|^2.
\]

\setcounter{theorem}{7}\begin{proposition}\label{prop:ergodicity-fourier}
In the setup above, a continuous surjective homomorphism $T\colon X\to X$
is $m_X$-invariant if and only if, for a character $\chi\in\widehat X$,
$\chi(T^n(x))=\chi(x)$ implies that $\chi$ is trivial, i.e. $\chi\equiv1$.
\end{proposition}
% Author check: the source says "m_X-invariant" rather than "ergodic" in
% this statement, and does not quantify n explicitly; retained without repair.
\begin{proof}
Suppose $T$ is ergodic and let $\chi$ be nontrivial such that
$\chi\circ T^n=\chi$, with $n$ the least such integer. Set
\[
  \widetilde\chi:=\chi+\chi\circ T+\cdots+\chi\circ T^{n-1}.
\]
Then $\widetilde\chi\circ T=\widetilde\chi$, i.e.
$U_T\widetilde\chi=\widetilde\chi$. From the operator-theoretic setting,
we see that $\widetilde\chi$ is constant. But
$\langle\widetilde\chi,\chi\rangle_X\neq0$, so $\widetilde\chi$ is not trivial.
Hence $\chi$ must be trivial.

Conversely, if $U_Tf=f$, first write $f\in L^2_{m_X}$ as
$f(x)=\sum_{\chi\in\widehat X}c_\chi\chi(x)$. Since $U_Tf=f$, we have
% Source page 11
$c_\chi=c_{\chi\circ T}=c_{\chi\circ T^2}=\cdots$.
Now, as $\sum_{\chi\in\widehat X}|c_\chi|^2$ converges to $\|f\|^2$, we must
have $c_\chi=0$, or $\{c_{\chi\circ T^n}\}_n$ is periodic. But that would
imply $\chi\equiv1$. Hence $f$ is a constant, so $T$ is ergodic.
\end{proof}
% Author check: the source's periodicity assertion concerns the coefficients
% c_{chi o T^n}, rather than explicitly the orbit of the characters.

\begin{remark}
Having seen ergodicity, it is natural to ask whether we can reduce a
measure-preserving system into ``ergodic components''. Given any $E\subset X$,
one has a natural $T$-invariant set $\{T^n(E)\}_{n\in\mathbb Z}$.
Now $X$ being ergodic assures us that $\{T^n(E)\}_{n\in\mathbb Z}$ covers $X$
almost everywhere. But if $X$ were not ergodic, there is no apparent reason
why a subset of $X$ would exist such that $\{T^n(F)\}_{n\in\mathbb Z}$ covers
$E$ almost everywhere, where $F\subseteq E$. If we hope to use some form of
Zorn's lemma, the existence of an ergodic subsystem
$(E,\mathcal B|_E,\mu|_E,T|_E)$ is required. Thus this problem of decomposing
measure-preserving systems into ergodic components is natural. We shall deal
with it later.
\end{remark}
% Author check: the source refers to the family {T^n(E)} as an invariant set,
% uses n in Z without assuming invertibility, and omits positive measure here.

\begin{remark}
Using Proposition~\ref{prop:ergodicity-fourier}---I mean the ideas in it---one
can prove that the circle-doubling map and the rotation map are ergodic.
Finally, as the Bernoulli system is isomorphic to the circle-doubling system,
the Bernoulli shift map is ergodic too. Thus almost all the examples we saw
are ergodic transformations.
\end{remark}
% Author check: no irrationality condition is given for "the rotation map"
% in this final remark; the wording has been retained.
% BLOG-CONTENT-END
\end{document}
